\section{An adjoint-parity obstruction}\label{sec:adjoint-parity}
The following elementary property of pivotal duality gives an additional
constraint on the signed tables. It does not require a braiding,
pseudo-unitarity or a choice of dimensions.

\begin{lemma}[Skew simples in adjoint products]\label{lem:adjoint-parity}
Let $\mathcal C$ be a pivotal fusion category over $\mathbb C$, and let
$X$ be a self-dual simple with $\nu_2(X)=-1$. For every object $Y$,
the multiplicity of $X$ in $Y^*\otimes Y$ is even.
\end{lemma}
\begin{proof}
Write $j_V:V\longrightarrow V^{**}$ for the pivotal structure. On
$\operatorname{Hom}(V,V^*)$ use the transpose $f^t=f^*j_V$.
The definition of the second indicator gives an isomorphism
$f:X\longrightarrow X^*$ with $f^t=-f$.
Put $W=Y\otimes X\otimes Y^*$ and
\[
 \alpha=j_Y\otimes f\otimes\mathrm{id}_{Y^*}:W\longrightarrow W^*.
\]
Duality reverses tensor factors. The identity
$j_Y^*j_{Y^*}=\mathrm{id}_{Y^*}$ and monoidality of $j$ therefore give
$\alpha^t=-\alpha$. Associativity and duality coherence maps are
understood in this formula.

On the complex vector space $\operatorname{Hom}(1,W)$ the pairing
\[
 B(v,w)=v^*\alpha w
\]
is nondegenerate: duality identifies this space with
$\operatorname{Hom}(W^*,1)$, composition with $\alpha$ is an
isomorphism, and the composition pairing on the unit isotypic component
is perfect by semisimplicity. Dualizing the scalar $B(v,w)$ and using
$v^{**}=j_Wv$ gives
$B(v,w)=w^*\alpha^t v=-B(w,v)$. Thus its dimension is even.
Rigidity identifies its dimension with
$\dim\operatorname{Hom}(X,Y^*\otimes Y)$, proving the assertion.
\end{proof}

\begin{corollary}[A congruence for sphericalization]\label{cor:signed-parity}
Let $R,T$ be an unsigned and signed pair as in
Proposition~\ref{prop:signed}, arising from a categorification.
For every self-dual basis element $b_i$ there is a sign
$\nu_i\in\{1,-1\}$ such that, for every $j$,
\begin{equation}\label{eq:signed-parity}
 N_{j,j^*}^{i}-\nu_i T_{j,j^*}^{i}\equiv0\pmod4.
\end{equation}
Consequently, if neither sign satisfies these congruences for some $i$,
the signed table cannot arise from a sphericalization.
\end{corollary}
\begin{proof}
The two lifts $X_i^+,X_i^-$ have opposite second indicators, as explained
in Section~\ref{sec:frontier-signed-gauss}. Take
$\nu_i=\nu_2(X_i^+)$. In the adjoint product of a chosen lift of $b_j$,
their multiplicities are respectively
$(N_{j,j^*}^{i}+T_{j,j^*}^{i})/2$ and
$(N_{j,j^*}^{i}-T_{j,j^*}^{i})/2$.
The lift whose indicator is $-1$ has even multiplicity by
Lemma~\ref{lem:adjoint-parity}; this is precisely
\eqref{eq:signed-parity}. The same $\nu_i$ must work for all $j$.
\end{proof}

The verifier applies this congruence to every self-dual basis element
of every retained signed table. A witness records an index for which
both indicator choices fail, and the checker recomputes the unsigned
and signed adjoint coefficients from the two tensors. Thus a parity
exclusion depends only on the fusion data. Its occurrences, including
the signed branches it resolves, are tabulated automatically in the
certificate ledger.

\subsection{Joint use of arithmetic, parity and moments}
The preceding criteria constrain the same hypothetical sphericalization.
They may therefore be imposed jointly, before expanding its possible
signed tables or indicator assignments.

\begin{corollary}[Joint necessary conditions for a signed lift]
\label{cor:joint-signed-conditions}
Let $R,T,\chi$ be as in Proposition~\ref{prop:signed}, and put
$K=\mathbb Q(\chi(b_i):i)$. If these data arise from a
sphericalization, then both of the following conditions hold.
\begin{enumerate}
\item Every formal codegree of $T$ is a $d$-number, and every value
of every character of $T$ is a cyclotomic integer.
\item The indicators in Corollary~\ref{cor:sign-certificate} satisfy
\[
 N_{j,j^*}^{i}-\nu_iT_{j,j^*}^{i}\equiv0\pmod4
 \qquad\text{for every }j
\]
for each self-dual nonunit $b_i$. One assignment of these indicators
and the twist signs $\eta_j$ must satisfy the signed-translate
moment inequalities at every real embedding of $K$.
\end{enumerate}
\end{corollary}
\begin{proof}
The characters of $T$ are the characters of the sphericalization's
Grothendieck ring on which the central sign object acts as $-1$.
Its formal codegrees on this summand are $2f$, with $f$ a formal
codegree of $T$. Proposition~\ref{prop:standard} applies:
multiplication by a nonzero rational number preserves the
$d$-number property, and the character values remain cyclotomic
integers.

The congruences are Corollary~\ref{cor:signed-parity} and concern
the same second indicators as the moment identities. Their choices
must therefore agree. Under Galois conjugation the indicators and
relevant twist signs remain fixed, since they lie in $\{1,-1\}$.
The dimensions remain real and the squared-dimensional measure
remains positive, so Proposition~\ref{prop:signed-translate-moment}
applies at every real embedding with that same assignment.
\end{proof}

\begin{corollary}[Parity-refined spherical bound]
\label{cor:parity-spherical-bound}
In Proposition~\ref{prop:spherical-codegree-moment}, replace $S$ by
\[
 S_{\rm par}=1+\sum_{i\in J}d_i+
 \sum_{\substack{i>0,\ i=i^*\\i\notin J}}|d_i|,
 \qquad
 J=\{i>0:i=i^*,\ N_{j,j^*}^{i}\text{ is odd for some }j\}.
\]
The resulting inequality remains necessary at every real embedding.
\end{corollary}
\begin{proof}
Lemma~\ref{lem:adjoint-parity} forces $\nu_2(X_i)=1$ for $i\in J$.
Thus $s=\sum_i d_i\nu_2(X_i)\leq S_{\rm par}$ at every real
embedding. The proof of Proposition~\ref{prop:spherical-codegree-moment}
uses this upper bound only after multiplication by the nonnegative
remaining mass $\Delta^2-h$, so the replacement preserves the
necessary inequality.
\end{proof}

\begin{corollary}[Parity-restricted signed enumeration]
\label{cor:parity-restricted-enumeration}
For a self-dual rank-four fusion ring, write the unsigned and signed
parameter tuples as $N=(N_0,\ldots,N_9)$ and $p=(p_0,\ldots,p_9)$
in the order $(a,b,c,d,e,f,g,h,i,j)$. Put
\[
 J_1=\{0,3,5\},\qquad J_2=\{1,6,8\},\qquad
 J_3=\{2,7,9\}.
\]
It suffices to search the union, over
$\nu=(\nu_1,\nu_2,\nu_3)\in\{1,-1\}^3$, of the domains
\[
 |p_k|\leq N_k,\qquad
 p_k\equiv\nu_iN_k\pmod4\quad(k\in J_i),\qquad
 p_4\equiv N_4\pmod2.
\]
One may intersect these domains with the prescribed odd-coefficient
positivity normalization. After duplicate and empty domains are
removed, the resulting at most eight boxes are pairwise disjoint.
\end{corollary}
\begin{proof}
Corollary~\ref{cor:signed-parity} supplies three indicator signs
for every categorical signed table. The sets $J_i$ are precisely
the coefficients of the three nonunit simples in the three adjoint
products. The remaining parameter $e$ retains its original signed
parity and size restrictions.

For fixed $i$, the domains for $\nu_i=1$ and $\nu_i=-1$
coincide if every $N_k$, $k\in J_i$, is even. Otherwise an odd
$N_k$ gives distinct residues modulo four, making these domains
disjoint. Intersecting with positivity restrictions preserves this
dichotomy, except that a domain may become empty. Since the sets
$J_i$ are disjoint, the assertion about the boxes follows.

A basis-sign change $s_i\in\{1,-1\}$ multiplies $T_{j,j}^{i}$
by $s_i$ and changes the chosen lift's indicator to $s_i\nu_i$.
The product $\nu_iT_{j,j}^{i}$ is unchanged. Thus the union is
preserved by basis-sign changes, and the positivity normalization
retains a representative of every categorical possibility.
\end{proof}
The omitted coefficient tuples violate adjoint parity, whether or not
they are associative. A search restricted to these domains therefore
counts parity-admissible normalized associative tables. Coincident
domains do not identify their indicator choices: all compatible
indicators remain available in the joint moment test. These refinements
use the same categorical criteria; they change their order of application
and the size of the finite searches.

The arithmetic conditions can also be applied once to a common bounded
set of signed tables before testing compatibility with any unsigned
ring. Appendix~\ref{sec:shared-signed-cover} proves completeness of
such a shared cover under basis permutations and sign changes. Every
compatible transform of a surviving orbit is still checked against
the unsigned ring's bounds, parities and remaining moment conditions.
